% 本文件由 LaTeX 排版 Agent 根据有效 translation.json 自动生成。
% author=Augustine of Hippo; work_id=1513; title=论斥责与恩宠

\chapter{论斥责与恩宠}

% structure: p0001 source=h1 target=omitted reason=page-title-already-rendered-by-parent

\Emph{由希波的圣奥思定写于公元426或427年}\par

\Emph{致瓦连廷及阿德卢梅图姆的隐修士们，成书一卷。}\par

% structure: p0004 source=h2 target=section reason=relative-heading-rank; base=h2

\section{奥思定订正录节录}

\Emph{取自奥思定\WorkTitle{订正录}第二卷第六十七章，关于以下论著\WorkTitle{论斥责与恩宠}。}\par

我又给同一批人写了另一篇论著，题为\WorkTitle{论斥责与恩宠}，因为有人告诉我，他们中有人说，对不遵守天主诫命的人不应加以斥责，而只应为他祈祷，使他能遵守。本书开头是：\Quote{我读了你们的信，至亲爱的瓦连廷兄弟。}\par

% structure: p0007 source=h2 target=section reason=relative-heading-rank; base=h2

\section{第一章——引言。}

我读了你的信——瓦连廷，我至亲爱的兄弟，以及与他一起事奉天主的你们——是你们的爱德通过弗洛鲁斯兄弟和与他同来的那些人送来的；我感谢天主，因为通过你们带给我的讲话，我得知你们在主内的平安、真理中的合一、以及爱中的热忱。但敌人曾在你们中试图颠覆一些人，靠着天主的仁慈和祂奇妙的良善，将敌人的诡计转为祂仆人的益处，反倒达成了这样的结果：你们中没有一人被击倒而变得更糟，反而有些人被建立得更好。因此，没有必要一再重新思考我已经在长篇论著中充分论证并传给你们的一切；因为你们的答复表明你们是如何接受的。尽管如此，绝不要以为读过一次就能充分熟悉它。因此，如果你们希望它极有成效，就不要因重读而感到厌烦，要通过反复阅读使之彻底熟悉；这样你们就能最准确地知道，那些问题的性质和种类是什么，为了解决并满足这些问题，有一个非人类的、而是神圣的权威出现，如果我们希望达到我们正趋向的目标，就不应偏离这个权威。\par

% structure: p0009 source=h2 target=section reason=relative-heading-rank; base=h2

\section{第二章——论法律、恩宠与自由意志的公教信仰。}

主自身不但指示我们应避免什么邪恶、应行什么善事——这是法律的字句所能成就的一切——而且还帮助我们避免邪恶、行善，这没有恩宠的圣神就无人能做到；如果缺少恩宠，法律就只会使我们有罪并杀死我们。因此宗徒说：\BibleQuote{文字叫人死，圣神却叫人活。}\BibleRef{格后 3:6}所以，合法地使用法律的人，在其中认识到善恶，并不依靠自己的力量，而是投奔恩宠，靠其帮助避免邪恶、行善。然而，除非\Quote{人的脚步由主安排，祂决定他的道路}，有谁会投奔恩宠呢？因此，渴望恩宠的帮助就是恩宠的开始；关于此，他说：\Quote{我说：现在我开始；这是至高者右手的改变。}所以要承认我们有行善行恶的自由选择；但在行恶时，每个人都脱离了正义，是罪的奴仆；而在行善时，除非被祂所释放，无人能自由；祂说：\BibleQuote{如果子使你们自由，你们就真是自由的了。}\BibleRef{若 8:36}这并非是说，当一个人从罪的辖制中被释放后，就不再需要祂救赎者的帮助；而是相反，他听从祂说：\BibleQuote{离了我，你们什么也不能做}\BibleRef{若 15:5}，他自己也对祂说：\Quote{作我的帮助者！不要离弃我。}我很高兴在我们兄弟弗洛鲁斯身上也发现了这信德，这无疑是真的、先知的和公教的信德；因此，那些误解他的人——我确实认为如今靠天主的恩宠他们已被纠正——更应该受到纠正。\par

% structure: p0011 source=h2 target=section reason=relative-heading-rank; base=h2

\section{第三章——什么是藉耶稣基督而来的天主的恩宠。}

因为必须领悟藉我们的主耶稣基督而来的天主的恩宠——即，只有藉此，人们才能脱离邪恶，没有它，他们就绝对做不了任何善事，无论是在思想、意志、情感或行为上；不仅是为了藉此恩宠的彰显，他们知道应做什么，而且是为了藉其助佑，他们以爱德做他们所知道的。确实，宗徒为那些他所说的人祈求善愿和善工的灵感：\BibleQuote{现在我们祈祷天主，使你们不作恶，不是要我们显得是经得起考验的，而是叫你们行善。}\BibleRef{格后 13:7}谁能听到这话而不醒悟并承认，我们离恶行善是来自主天主呢？——因为宗徒确实不是说：我们劝告、教导、鼓励、斥责；而是说：\BibleQuote{我们祈祷天主，使你们不作恶，而是叫你们行善。}\BibleRef{格后 13:7}然而他也惯常向他们说话，做我提到过的所有那些事——他劝告、教导、鼓励、斥责。但他知道，他公开栽种、浇灌的一切，除非那位暗中赐予生长的天主应允他为他们的祷告，否则毫无用处。因为，正如这位外邦人的老师所说：\BibleQuote{可见栽种的算不得什么，浇灌的也算不得什么，只在那使之生长的天主。}\BibleRef{格前 3:7}\par

% structure: p0013 source=h2 target=section reason=relative-heading-rank; base=h2

\section{第四章——天主的儿女由天主的圣神引导。}

因此，愿那些问\Quote{为何给我们宣讲和命令要弃恶行善，如果这不是我们自己做的，而是\InnerQuote{天主在我们内工作，使我们愿意并实行}？}的人不要自欺。\BibleRef{斐 2:13} 但让他们明白，如果他们是天主的子女，就被天主的神引导\BibleRef{罗 8:14}，去做当做的事；当他们做了，就当感谢使他们行动的那一位。因为他们被推动而行动，并非要他们自己无所作为；此外，也指示他们当做的事，好使他们按当行的去行——即怀着对正义的爱与喜悦——并因领受了\Quote{主所赐的甘甜，使他们的土地出产果实}而欢喜。但当他们不行动时，或因完全不行动，或因不怀着爱行动，就让他们祈祷，好获得尚未拥有的。因为他们将有什么不是领受的呢？或者他们有什么不是领受的呢？\BibleRef{格前 4:7}\par

% structure: p0015 source=h2 target=section reason=relative-heading-rank; base=h2

\section{第五章 [III.]——劝责不可疏忽。}

\Quote{那么，}他们说，\Quote{愿那些管理我们的人只命令我们当做的事，并为我们祈祷，好使我们能做；但若我们不做，就不要劝责和谴责我们。}当然一切都当实行，因为教会教师们，即宗徒们，惯于做一切事——既命令当做的事，若未做则劝责，并祈祷使能做成。宗徒命令说：\Quote{你们一切事都要以爱德而行。}\BibleRef{格前 16:14} 他劝责说：\Quote{所以，你们彼此有诉讼的事，已是全然缺失。为什么不情愿受欺？为什么不情愿吃亏？你们反倒欺压和亏负，且是亏负弟兄。你们岂不知不义的人不能承受天主的国吗？}让我们也听他祈祷：\Quote{愿主}，他说，\Quote{使你们彼此相爱的爱德和对众人的爱德增长并满溢。}\BibleRef{得前 3:12} 他命令应持守爱德；他劝责，因为爱德未持守；他祈祷，使爱德丰盈。人啊！借着他的命令学习你应当拥有什么；借着他的劝责学习你没有它乃是因你自己的过失；借着他的祈祷学习你从哪里可领受你所渴望拥有的。\par

% structure: p0017 source=h2 target=section reason=relative-heading-rank; base=h2

\section{第六章 [IV.]——反对使用劝责的异议。}

\Quote{怎么，}他说，\Quote{我没有从祂领受的，怎么是我的过失？因为除非由祂赐予，否则这样大的恩赐从何处能有？}请允许我稍作争辩，我的弟兄们，不是反对你们——你们的心在天主前是正直的，而是反对那些思念世俗之事的人，或反对那些人的思考方式本身，为维护属天和神圣的恩宠。因为说这话的人，是在他们的恶行中不愿意被那些宣讲这恩宠的人劝责。\Quote{命令我当做的事，我若做了，就为我感谢天主，因祂使我做了；我若不做，就不应被劝责，而应恳求祂赐予祂所未赐的，即那遵守祂诫命的、对天主和近人的信德之爱。那么，为我祈祷，使我领受它，并借此自由而甘愿地遵行祂所命令的。但若因我自己的过失而没有它，我才应受公正的劝责；即若我能自己给自己，或能领受却没有做，或是祂赐予而我却不愿领受。既然连意志本身也是由主预备的，\BibleRef{箴 16:1} 你为什么因见我不愿遵行祂的诫命而劝责我，而不亲自求祂也在我内推动这意愿呢？}\par

% structure: p0019 source=h2 target=section reason=relative-heading-rank; base=h2

\section{第七章 [V.]——劝责的必要性与益处。}

我们回答：无论你是谁，你不遵行已知道的天主诫命，且不愿受劝责，你正因此——因你不愿受劝责——而必须被劝责。因为你不愿你的过犯被指出；你不愿它们被触及，不愿这种有益的痛苦加于你，使你去寻求医生；你不愿被显示给自己，使你在看见自己丑恶时，渴望改新者，并恳求祂使你不继续处于那种可恶。因为你是恶的，这是你的过犯；更大的过犯是，因你是恶而不愿受劝责，仿佛过犯应当被称赞，或被漠视既不称赞也不责备，又仿佛被劝责者的惧怕、羞耻或痛悔毫无益处，或除了健康地激发人恳求那善者，以致被劝责的恶人变成可称赞的善人之外，还有其他益处。因为那个不愿受劝责的人希望为他做的事，如他说：\Quote{不如为我祈祷}——他正因此必须被劝责，使他也能为自己去做；因为那种痛悔——当他感到劝责的刺时对自己不满意——会激发他渴望更恳切的祈祷，好借天主的仁慈，因爱德的增长获得帮助，停止做可耻和令人痛悔的事，而做可称赞和令人喜悦的事。这就是有益施行的劝责之益处，有时按罪的多样性而施以更大或更小的严厉；并且它是有益的，当至高医生垂顾时。因为除非它使人痛悔己罪，否则毫无益处。谁赐予这个呢，岂不是那曾垂顾否认祂的宗徒伯多禄，使他痛哭的那一位？\BibleRef{路 22:61} 因此，宗徒保禄在说对那些持不同见解的人应以温和劝责之后，立即加上：\Quote{但愿天主赐予他们悔改，得以认识真理，并挣脱魔鬼的罗网。}\BibleRef{弟后 2:25}\par

% structure: p0021 source=h2 target=section reason=relative-heading-rank; base=h2

\section{第八章——进一步答复那些反对劝责的人。}

但是，那些不愿受责斥的人为何说：\Quote{尽管吩咐我，并为我祈祷，使我能够遵行你的吩咐}？他们为何不本着邪恶倾向，也同样拒绝这些，说：\Quote{我愿你既不要吩咐我，也不要为我祈祷}？有谁曾为伯多禄祈祷，求天主赐他悔改，使他哀哭自己否认主的罪？有谁教训保禄关于属于基督信仰的神圣诫命？因此，当他被人听见宣讲福音并说：\Quote{弟兄们，我告诉你们，我所宣讲的福音不是出于人。因为我既不是从人领受的，也不是从人学来的，而是藉耶稣基督的启示来的，} \BibleRef{迦 1:11} —— 难道会有人回答他说：\Quote{你为何麻烦我们，要我们从你领受并学习那些你既非从人领受也未从人学来的东西？那赐给你的，也能照样赐给我们}？此外，如果他们不敢这样说，而容许福音藉人传给他们，尽管它不能由人传给另一个人，那么，让他们也承认，他们应当受那些管辖他们的人责斥，这些人宣讲基督的恩宠；尽管不能否认，即使无人责斥，天主也能纠正祂所愿意的人，并以祂药物最隐秘而强大的力量引导他达到有益的悔改克苦。正如我们不应停止为那些我们希望被纠正的人祈祷——尽管没有任何人为伯多禄祈祷，主却注视他，使他哀哭自己的罪——同样，我们也不可忽略责斥，尽管天主能使祂所愿意的人甚至在不受责斥时就被纠正。但当那怜悯并帮助的天主，使祂所愿意的人甚至在无责斥时也受益时，人就藉责斥受益。至于为何这些人按一种方式、那些人按另一种方式、其他人又按另一种方式被召叫归正，方式各异、数不胜数，我们切不可断定这是泥土的事，而是陶匠的事。\par

% structure: p0023 source=h2 target=section reason=relative-heading-rank; base=h2

\section{第九章 [VI]——为何那些尚未领受服从恩宠而不服从天主的人，应被公正地责斥。}

他们说：\Quote{宗徒说：\InnerQuote{有什么使你与人不同呢？你有什么不是领受的呢？既然是领受的，为何你夸耀，好像不是领受的呢？} \BibleRef{格后 4:7} 那么，为何我们要受责斥、指责、训斥、控告？我们这些没有领受的人做了什么？}说这话的人想显得自己违背天主是无辜的，因为服从本身确实是祂的恩赐；这恩赐必然存在于拥有爱德的人内，而这爱德无疑是出于天主，\BibleRef{若壹 4:7}，且圣父将之赐给祂的子女。他们说：\Quote{这个我们没有领受。那么，为何我们受责斥，好像我们能给自己这恩赐，却凭自己的选择不给？}他们却没有注意到，如果他们尚未重生，那么当他们因违背天主而受责斥时，首先应当自责的是：天主从人类受造之初就使人正直，\BibleRef{训 7:30}，而天主并没有不义。\BibleRef{罗 9:14}因此，最初的败坏——即不服从天主——是出于人，因为人凭自己的邪恶意志从起初天主的正直中堕落，就成了败坏。难道这败坏在一个人身上就不该被责斥，因为它并非仅属受责者，而是人人共有吗？不，即使在个人身上，那人人共有的也该被责斥。因为没有人完全摆脱它，这并不妨碍它附着于每个人。那些原罪确实被称为他人的罪，因为各人从父母继承了它们；但它们也被合理地说成是我们自己的，因为如宗徒所说，在那一人内众人犯了罪。\BibleRef{罗 3:23}因此，让那该受罚的根源受责斥，以便从责斥的痛悔中生出重生的意愿——如果那受责斥的人确实是恩许之子，这样，藉责斥从外发出的声音和鞭打，天主能从内在藉隐藏的默感也在他内激发意愿。但是，如果一个人已经重生而成义，却凭自己的自由选择堕落于恶生，他肯定不能说\Quote{我没有领受}，因为他凭自己的自由选择作恶，他已失去了他已领受的天主恩宠。如果他因责斥而痛悔并有益地哀哭，回到相似甚至更善的善工，那么责斥的益处在这里就极其明显。然而，藉着人之手实行的责斥能否生效，无论是出于爱德与否，唯靠天主。\par

% structure: p0025 source=h2 target=section reason=relative-heading-rank; base=h2

\section{第十章——一切坚忍都是天主的恩赐。}

像这样不愿受劝责的人，还能说：\Quote{我没有领受什么，我做了什么？}尽管显然他已领受，并且因自己的过失失去了他所领受的？他说：\Quote{我能，我完全能——当你责备我因自己的意志从善生堕落为恶生时——我仍能说：\InnerQuote{我没有领受什么，我做了什么？}因为我领受了藉爱德行动的信德，但没有领受持守至终的恩赐。有谁敢说这持守不是天主的恩赐，并且如此伟大的产业是我们在这样意义上拥有的，以至于若有人拥有它，宗徒就不能对他说：\InnerQuote{你有什么不是领受的呢？}\BibleRef{格前 4:7}既然他以这样的方式拥有它，好像他没有领受它？}对此，我们确实无法否认：在善中持守，甚至进步至终，也是天主的一大恩赐；并且它只来自祂，正如经上所载：\BibleQuote{一切美好的赠与、一切完善的恩赐，都是从上而来，从光明之父降下的。}\BibleRef{雅 1:17}但不可因此忽略对未能持守者的劝责，\BibleQuote{免得天主或许赐予他悔改，使他脱离魔鬼的罗网；}\BibleRef{弟后 2:25}因为劝责的益处，宗徒附加了这判断，如我上文所说：\BibleQuote{以温和劝责那些思想不同的人，免得天主有时赐予他们悔改。}\BibleRef{弟后 2:25}因为若我们说，如此可赞美、如此有福的持守，是在人自己身上而不是从天主而来，那么首先，主对\Person{伯多禄}所说的话就落空了：\BibleQuote{我已为你祈求，为叫你的信德不致丧失。}\BibleRef{路 22:32}祂为他求的是什么，岂非持守至终吗？确实，如果人能靠自己获得这恩赐，就不必向天主祈求了。然后，当宗徒说：\BibleQuote{现在我们祈祷天主，使你们不作恶事，}\BibleRef{格后 13:7}毫无疑问，他是为他们向天主祈求持守。因为那离弃善、不持守于善而转向恶的人，确实\BibleQuote{不作恶事}？再者，那里他说：\BibleQuote{我每逢想念你们，就感谢我的天主，在每一次祈祷中，常为你们众人以喜乐祈求，因为你们从最初直到现在，在福音上共融；我深信那在你们内开始美好工作的，必完成它直到耶稣基督的日子，}——这不是从天主的仁慈向他们许诺持守至终，又是什么？又那里他说：\BibleQuote{\Person{厄帕夫辣}问候你们，他是你们中的一位，基督耶稣的仆人，常在祈祷中为你们竭力，使你们在一切天主的旨意中得以圆满成就而站立，}\BibleRef{哥 4:12}——\BibleQuote{使你们站立}不是\BibleQuote{使你们持守}是什么？因此论到魔鬼曾说：\BibleQuote{他在真理中不站立，}\BibleRef{若 8:24}因为他曾在那里，却没有持续。因为那些人确实已经站在信德中。当我们为站立的人祈祷使他站立，我们不是祈求别的，而是祈求他持守。\Person{犹达}宗徒又说：\BibleQuote{但愿那能保护你们不失足，并使你们无瑕无疵地立在他荣耀面前，在喜乐中，}\BibleRef{犹 1:24}——这不是最清楚表明持守至终是天主的恩赐吗？因为那保护人不失足，以使他无瑕无疵地立在荣耀面前、在喜乐中的，除了持守于善，还赐什么？再者，我们在宗徒大事录中读到：\BibleQuote{外邦人听见了，就喜欢，领受了主的话；凡预定得永生的人都信了。}\BibleRef{宗 13:48}除了持守的恩赐，谁能预定得永生？当我们读到：\BibleQuote{唯独持守到底的，才可得救；}\BibleRef{玛 10:22}那是何等救恩，岂非永生？当我们在主祷文中对天主父说：\BibleQuote{愿你的名被尊为圣，}\BibleRef{玛 6:9}我们求的是什么，岂非祂的名在我们内被尊为圣？既然这事已借着重生之洗完成，为何信徒每日仍求？除非是使我们持守于那在我们内已成就的事。因为真福\Person{西彼廉}在解释同一祈祷时也如此理解，他说：\Quote{我们说\InnerQuote{愿你的名被尊为圣}，并非我们愿天主因我们的祈祷而成圣，而是我们向天主祈求，使祂的名在我们内被尊为圣。但天主由谁成圣呢？因祂自己成圣？是的，因祂曾说：\InnerQuote{你们应当是圣的，因为我是圣的。}我们祈求恳求，使我们已在圣洗中成圣的，能持守于我们已开始成为的。}看，这位最光荣的殉道者如此认为：基督的信徒每日以这些话所求的，就是持守于他们已开始成为的。无人需要怀疑，凡向主祈求持守于善的，就承认这持守是祂的恩赐。\par

% structure: p0027 source=h2 target=section reason=relative-heading-rank; base=h2

\section{第十一章 [VII.]——凡未领受持守恩赐，且已陷于大罪并死于其中者，理当受义罚。}

如果事情是这样，我们仍然劝诫那些人，并且合理地劝诫他们，他们虽然曾经生活良好，却没有在其中恒心持守；因为他们出于自己的意志从善的生活改变为恶的生活，因此应受劝诫；而且如果劝诫对他们毫无益处，他们持续在毁灭的生活中直到死亡，他们也应受永远的、神圣的定罪。他们也不能以这样的借口为自己辩护，说——正如他们现在所说：\Quote{为何我们被劝诫？}——那么，\Quote{为何我们被定罪，既然实际上，我们之所以从善回到恶，是因为我们没有领受那使我们能留在善中的恒心？}他们绝不能藉此借口逃脱公义的定罪。因为如果按照真理之言，没有人能逃脱因亚当而招致的定罪，除非通过耶稣基督的信德，然而那些能说他们没有听过基督福音的人，也不能从这定罪中逃脱，理由是\BibleQuote{信德出于听闻}\BibleRef{罗 10:17}，更何况那些说\Quote{我们没有领受恒心}的人，又怎能逃脱呢？因为那些说\Quote{我们没有领受听闻}之人的借口，似乎比那些说\Quote{我们没有领受恒心}之人的借口更为公平；因为可以对一个人说：在你所听见并持守的事上，你若愿意，你本可以恒心持守；但绝不能对一个人说：你所未曾听见的，你若愿意，你本可以相信。\par

% structure: p0029 source=h2 target=section reason=relative-heading-rank; base=h2

\section{第十二章 那些没有领受恒心的人与丧亡的群体并无区别}

因此，无论是没有听过福音的人，还是听了福音并因而改善、却没有领受恒心的人，以及听了福音却拒绝来到基督面前（即相信祂），因为祂亲自说：\BibleQuote{除非蒙我父赐给他，没有人能到我这里来}\BibleRef{若 6:65}，还有那些因年幼不能相信、但本可藉重生之洗解除原罪、却没有领受这洗而死于死亡中的人：所有这些人，都没有与那显然被定罪的同一团块区别开来，因为所有人都从一人进入定罪。然而，有些人得以区别开来，不是靠他们自己的功劳，而是靠中保的恩宠；也就是说，他们是在第二亚当的血中白白成义的。因此，当我们听到：\BibleQuote{谁使你与众不同？你有什么不是领受的呢？既然是领受的，为什么你夸耀，好像不是领受的一样？}\BibleRef{格前 4:7}时，我们应该明白：从那源于第一亚当的丧亡群体中，没有人能与众不同，除非他拥有这份恩赐，而凡拥有这恩赐的，都是因救主的恩宠领受的。这一见证非常重要，以致真福西彼廉在写给奎里努斯的信中将其作为标题，他说：\Quote{我们不可在任何事上夸耀，因为没有任何事物是我们自己的。}\par

% structure: p0031 source=h2 target=section reason=relative-heading-rank; base=h2

\section{第十三章 拣选是出于恩宠，而非出于功劳}

因此，凡是从那原始的定罪中藉着如此丰厚的神圣恩宠而得以区别出来的人，毫无疑问，天主为他们提供了这样的安排：他们应当听到福音，当他们听到时便相信，并在那藉着爱德运行的信德中恒心持守到底；而且，如果他们偶尔偏离了道路，当他们受到劝诫时，他们得以改正，其中有些人，即使没有受到人的劝诫，也会返回他们曾离开的道路；还有一些在任何年龄领受了恩宠的人，因速死的缘故被从这个生命的危险中撤走。因为这一切都是祂在他们内施行的，祂使他们成为仁慈的器皿，祂也在创世之前，在祂的圣子内，藉着恩宠的拣选拣选了他们：\BibleQuote{既是出于恩宠，就不再是由于行为，不然，恩宠就不成其为恩宠了。}\BibleRef{罗 11:6}因为他们并不是如此蒙召以致不被拣选，关于此有话说：\BibleQuote{因为被召的人多，被选的人少。}\BibleRef{玛 20:16}但因他们按照旨意蒙召，他们必定也是藉着恩宠的拣选而被拣选，正如经上所说，不是出于他们任何先前的功劳，因为对他们而言，恩宠就是一切功劳。\par

% structure: p0033 source=h2 target=section reason=relative-heading-rank; base=h2

\section{第十四章 蒙拣选和预定的人无一可以丧亡}

宗徒论及这样的人说：\BibleQuote{我们知道，凡爱天主的，祂使万事互相效力，使之得益，就是那些按祂旨意蒙召的人。因为祂所预知的人，也预定他们与祂圣子的肖像相似，使祂在众多弟兄中作长子。此外，祂所预定的人，也召叫了他们；祂所召叫的人，也使他们成义；祂所成义的人，也使他们受光荣。}这些人没有一个丧亡，因为都是蒙拣选的。他们之所以蒙拣选，是因为他们按旨意蒙召——这旨意不是他们的，而是天主的。关于这旨意，祂在别处说：\BibleQuote{为使天主拣选的旨意坚定，不因人的作为，而因那召叫人的，就对她说过：大的要服事小的。}\BibleRef{罗 9:11} 又在另一处说：\BibleQuote{不是按我们的作为，而是按祂自己的旨意和恩宠。}\BibleRef{弟后 1:9} 因此，当我们听到\BibleQuote{此外，祂所预定的人，也召叫了他们，}\BibleRef{罗 8:29} 我们应当承认他们是按祂的旨意蒙召的；因为祂从那处开始说：\BibleQuote{祂使万事互相效力，使那些按祂旨意蒙召的人得益，}然后又加上：\BibleQuote{因为祂所预知的人，也预定他们与祂圣子的肖像相似，使祂在众多弟兄中作长子。}对这些许诺，祂又加上：\BibleQuote{此外，祂所预定的人，也召叫了他们。}因此，祂希望这些人是按祂旨意蒙召的人，免得他们中有人因主的那句话而以为自己是蒙召而未蒙拣选的：\BibleQuote{被召的人多，蒙拣选的人少。}\BibleRef{玛 20:16} 因为凡蒙拣选的，毫无疑问也是蒙召的；但并非所有蒙召的都必然蒙拣选。因此，如常言所说，那些按旨意蒙召、且被预定和预知的人，才是蒙拣选的。如果这些人中有一个丧亡，天主就犯了错误；但他们中没有一个丧亡，因为天主不会犯错。如果这些人中有一个丧亡，天主就被人的罪战胜了；但他们中没有一个丧亡，因为天主不会被任何事物战胜。此外，他们蒙拣选是要与基督一同为王，不像犹达斯那样，他被选去执行一项适合他的任务。因为他是被那一位所选的，祂深知如何利用恶人，甚至通过他可诅咒的行为，完成那为了祂自己而来的可敬工程。因此，当我们听到\BibleQuote{我不是拣选了你们十二个人吗？你们中间却有一个是魔鬼，}\BibleRef{若 6:70} 我们应当明白，其余的人因仁慈而蒙拣选，而他因审判而蒙拣选；那些人是为了获得祂的王国，他是为了流祂的血！\par

% structure: p0035 source=h2 target=section reason=relative-heading-rank; base=h2

\section{第十五章——恒心赐予到底。}

接下来，向蒙拣选者国度所说的话是贴切的：\BibleQuote{如果天主偕同我们，谁能反对我们？祂既然没有怜惜自己的圣子，反而为我们众人把祂交出来，岂不也把一切与祂一同赐给我们？谁能控告天主所拣选的人？是天主使他们成义的吗？谁能定他们的罪？是那已死，而且更已复活，如今在天主右边，代我们转求的基督耶稣吗？} 他们领受了何等坚固、甚至到末了的恒心的恩赐，就让他们接着说：\BibleQuote{谁能使我们与基督的爱隔绝？是困苦吗？是窘迫吗？是迫害吗？是饥饿吗？是赤贫吗？是危险吗？是刀剑吗？正如经上所载：\InnerQuote{为了你，我们终日被杀害，被视为待宰的羊。}然而，靠那爱我们的主，我们在这一切事上大获全胜。因为我深信，无论是死，是生，是天使，是掌权者，是现在的事，是将来的事，是权势，是崇高，是深远，或是任何别的受造物，都不能使我们与天主的爱隔绝，这爱是在我们的主基督耶稣内的。}\par

% structure: p0037 source=h2 target=section reason=relative-heading-rank; base=h2

\section{第十六章——凡不恒心到底的人，并不因预定而从灭亡的群体中被分别出来。}

这样的人就是那些在写给弟茂德的书信中所指的。那里说：赫默纳约和斐肋托颠覆了一些人的信德，随即加上：\BibleQuote{然而，天主的坚固基础屹立不摇，带着这样的印记：主认识属于祂的人。}\BibleRef{弟后 2:19} 这些人的信德，以爱德行事，要么根本不会动摇；即使有些人的信德动摇了，也会在生命终结之前恢复，已介入的罪孽被除去，并且获得恒心至终的恩赐。但那些不将持守到底，且将如此远离基督徒的信德和行为，以致生命的终结发现他们处于那种状态的人，毫无疑问，即使在那个他们生活得好而虔诚的时刻，也不可算在这些人的数目中。因为借着天主的预知和预定，他们并没有从灭亡的群体中被分别出来，因此不是按天主的旨意蒙召，所以不是蒙拣选的；而是在那些被称为\BibleQuote{被召的人多}的人之中，而不是在那些被称为\BibleQuote{但蒙拣选的人少}的人之中。然而，谁能否认他们是蒙拣选的，既然他们相信、受洗，并按照天主生活？显然，他们被那些不知道他们将来如何的人称为蒙拣选的，但并非被那一位称为蒙拣选的，那一位知道他们不会拥有带领蒙拣选者进入永福的恒心，并且知道他们站着，正如祂预知他们会跌倒一样。\par

% structure: p0039 source=h2 target=section reason=relative-heading-rank; base=h2

\section{第十七章——为什么恒心赐予一人而不赐予另一人是不可探究的。}

此处，若有人问我：为何天主没有将\OriginalTerm{持守}{perseverantia}之恩赐给那些祂曾赐予爱、使他们能按基督徒方式生活的人？我答说：我不知道。因我并非妄自尊大，而是承认自己度量渺小地听到宗徒说：\BibleQuote{人啊！你是谁，竟敢向天主抗辩？}\BibleRef{罗 9:20} 以及：\BibleQuote{啊！天主的智慧和知识的富藏，是多么高深！他的判断是多么不可测量，他的道路是多么不可追踪！}\BibleRef{罗 11:33} 因此，就祂屈尊向我们显示祂的判断而言，让我们感恩；就祂认为适宜隐藏它们而言，让我们不要抱怨祂的旨意，但要相信这对我们也是最为有益的。然而，无论你是谁，若你敌视祂的恩宠而这样质问，你自己要说什么？好在你不否认自己是基督徒，并自夸为公教信徒。因此，若你承认持守至终于善是天主的恩赐，我想你同我一样，不知为何一人应接受此恩赐而另一人不应接受；在此情形下，我们俩都无法窥测天主不可追踪的判断。或者，若你说持守于善或不持守属于人的\OriginalTerm{自由意志}{liberum arbitrium}——你护卫这自由意志，不是合乎天主的恩宠，而是反对它——并且若人持守并非由于天主的恩赐，而是由于人意志的作为，那么你为何要反对那说这话者的言语：\BibleQuote{我已为你祈求了，伯多禄，为叫你的信德不致丧失}\BibleRef{路 22:32}？你敢说，即使基督祈求伯多禄的信德不致丧失，若伯多禄愿意它丧失，它仍会丧失吗？即若他不愿它持续至终？好像伯多禄能丝毫愿意不同于基督为他所求使他愿意的。因为谁不知道，若那使他忠诚的意志丧失，伯多禄的信德便会灭亡；若那同一意志存留，信德就会存留？但因\BibleQuote{人的意志由上主所预备}\BibleRef{箴 8:35}，因此基督为他所作的祈求不能是虚妄的祈求。那么，当祂祈求他的信德不致丧失时，祂所求的是什么？岂非就是在信德中他拥有最自由、最坚强、不可战胜、持守不渝的意志吗？看，自由意志在何等的程度上是合乎天主的恩宠而非反对它而被护卫；因为人的意志不是藉自由获得恩宠，而是藉恩宠获得自由，并获得令人喜悦的恒心和不可战胜的刚毅，以便能持守。\par

% structure: p0041 source=h2 target=section reason=relative-heading-rank; base=h2

\section{第十八章——天主奇妙审判的一些实例}

确实，这令人惊奇，大大令人惊奇：对于一些祂自己的子女——就是祂在基督内所重生的，祂曾赐予信、望、爱的人——天主也没有赐予持守，而对祂人的子女，祂竟赦免如此邪恶，并藉施予恩宠使他们成为祂自己的子女。谁不对此感到惊奇？谁不因此极其震惊？此外，同样奇妙、真实且显而易见，甚至连天主的恩宠的仇敌也无法否认的是：祂的一些朋友（即重生且善良的信徒）的子女，在婴儿时期未受圣洗而离开此世——虽然祂当然可以恩赐这圣洗之恩，如果祂愿意，因为一切都在祂权能之下——祂使他们不能进入祂的国度，而祂却引入了他们的父母；祂的一些仇敌的子女，祂使他们落入基督徒手中，并藉此圣洗之恩引入国度，而他们的父母却是局外人；虽然对于前一种婴儿，他们并无任何本身意志所招致的恶；对于后一种，亦无任何本身意志所成就的善。诚然，在此情形中，天主的判断，因其公义且深奥，既不可责备，亦不可揣测。其中也包括我们现在所讨论的关于持守的判断。因此，关于两者，我们可以高呼：\BibleQuote{啊！天主的智慧和知识的富藏，是多么高深！他的判断是多么不可测量！}\BibleRef{罗 11:33}\par

% structure: p0043 source=h2 target=section reason=relative-heading-rank; base=h2

\section{第十九章——天主之道路不可测度}

也不要因我们无法测度祂不可探究的途径而惊讶。因为，且不提主天主赐给一些人而未赐给其他人的无数其他事物——在祂那里没有偏袒，这些事物并非凭意志的功德而赐予，如身体的敏捷、力量、健康、美貌，令人惊奇的理智和能掌握许多技艺的心智禀赋，或从外部降临到人身上的事物，如财富、高贵、荣誉及其他此类，只有天主的能力才能使人拥有；甚至也不提婴儿的圣洗圣事（那些反对者中没有一人能说，如同可以说其他事物那样，这不属于天主的国）——为何赐给这个婴儿而不赐给那个，既然两者都同样在天主的能力中，且没有这圣事无人能进入天主的国——那么，暂且不提这些事，或把它们搁置一边，让人思考我们正在讨论的那些非常特殊的情况吧。因为我们正在谈论的是那些没有坚持行善到底，而是在善意的衰落中从善堕入恶而死去的人。让反对者回答，如果可能的话，为什么当这些人虔诚地生活时，天主没有在那时将她们从这生命的危险中夺去，\BibleQuote{免得邪恶改变了她们的心意，免得诡诈迷惑了她们的灵魂}？\BibleRef{智 4:11}难道祂没有能力这样做，还是祂不知道她们将来的罪恶？的确，除非极其乖谬和疯狂，没有人会这样说。那么，为什么祂没有这样做呢？让那些嘲笑我们的人回答吧，当我们在此类事上高呼：\BibleQuote{祂的判断是多么不可测度，祂的道路是多么难以寻觅！}\BibleRef{罗 11:33}因为要么天主将此赐予祂所愿意赐予的人，要么那圣经无疑是错误的，它论及义人的早逝时说：\BibleQuote{他已被夺去，免得邪恶改变了他的心意，或者免得诡诈迷惑了他的灵魂。}\BibleRef{智 4:11}那么，为什么天主将如此巨大的恩惠赐给一些人，而不赐给另一些人呢？须知在祂那里没有不义\BibleRef{罗 9:14}，也没有偏袒人\BibleRef{罗 2:11}，并且祂有能力决定每个人在这被称为世上考验的生命中停留多久？\BibleRef{约 7:1}因此，既然他们被迫承认，一个人在从善变为恶之前结束此生是天主的恩赐，但他们不知道为何赐给一些人而不赐给另一些人，那么，让他们与我们一同承认，坚持行善是天主的恩赐，按照圣经——我已从中引用了许多见证——并且让他们谦卑地与我们一同，在不知为何赐给一些人而不赐给另一些人的事上，不对天主发怨言。\par

% structure: p0045 source=h2 target=section reason=relative-heading-rank; base=h2

\section{第二十章 [IX.]——有些人是按照暂时领受的恩宠成为天主的儿女，有些人是按照天主永恒的预知。}

也不要因天主没有将这恒心赐给祂的一些儿女而困扰我们。然而，如果他们是那些按预定蒙召、按祂旨意蒙召的人——他们确实是恩许的儿女——就绝不会如此。因为前者，当他们虔诚生活时，被称为天主的儿女；但由于他们将邪恶地生活，并死于那种不虔敬中，天主的预知并不称他们为天主的儿女。因为有些人是天主的儿女，我们尚未拥有，天主却已拥有，关于他们，福音作者若望说：\BibleQuote{耶稣将为那民族而死，并且不只为那民族，还要将天主四散的儿女聚集归一。}\BibleRef{若 11:51}这当然是他们要通过相信、借着福音的宣讲而成为的。然而在这发生之前，他们已经在天父的纪念中，以不可改变的坚定被登记为天主的儿子。再者，有些人因着甚至在暂时事物上领受的恩宠，被我们称为天主的儿女，却不被天主如此称呼；关于他们，同一位若望说：\BibleQuote{他们从我们中间出去，但却不是属于我们的；因为如果他们属于我们，就必定会与我们一同存留。}\BibleRef{若一 2:19}他没有说：\Quote{他们从我们中间出去，但因他们未与我们一同存留，现在就不再属于我们了}；而是说：\BibleQuote{他们从我们中间出去，但却不是属于我们的}——也就是说，即使当他们出现在我们中间时，他们也不属于我们。并且好像有人对他说：你如何证明这一点？他说：\BibleQuote{因为如果他们属于我们，就必定会与我们一同存留。}\BibleRef{罗 8:29}这是天主儿女的话；是若望在说话，他是在天主儿女中被指定居首位的人。因此，当天主的儿女论到那些没有恒心的人说：\BibleQuote{他们从我们中间出去，但却不是属于我们的}，并加上：\BibleQuote{因为如果他们属于我们，就必定会与我们一同存留}，这难道不是说，即使当他们处于儿女的宣信和名号时，他们也不是儿女吗？不是因为他们假装正义，而是因为他们没有在正义中存留。因为若望并没有说：\Quote{因为他们如果属于我们，就必定会与我们一同保持真实而非伪装的正义}；而是说：\BibleQuote{如果他们属于我们，就必定会与我们一同存留。}毫无疑问，他希望他们在善中存留。因此，他们曾处于善中；但因他们没有在善中存留——也就是说，他们没有坚持到底——他说，当他们与我们同在时，他们并不属于我们——也就是说，即使当他们处于儿女的信德时，他们也不在儿女的行列中；因为真正的儿女，是预知和预定要肖似祂儿子的肖像，且按祂旨意蒙召，以致被拣选的人。因为恩许之子不会丧亡，而丧亡之子才会丧亡。\BibleRef{若 17:12}\par

% structure: p0047 source=h2 target=section reason=relative-heading-rank; base=h2

\section{第二十一章——谁可以被理解为赐予基督的人。}

那么，这些人属于蒙召的群众，却不属于少数被选的人。因此，天主并没有将恒心赐给祂所预定的子女，因为如果他们是在那为数之中的子女，他们就必会拥有它；而他们若没有领受，又怎能拥有呢？正如经上说：\BibleQuote{你有什么不是领受的呢？}（\BibleRef{格前 4:7}）因此，这样的子女被赐予基督圣子，正如祂亲自对圣父说：\BibleQuote{祢赐给我的人，愿他们一个也不失丧，反而获得永生。}（\BibleRef{玛 20:16}）所以，那些被预定得永生的人，被理解为是赐予基督的。这些人就是按旨意蒙召、被预定的，其中没有一人丧亡。因此，当他们从善变为恶时，他们不会以此终结此生，因为他们是这样被预定、为此目的被赐予基督的，好使他们不失丧，反而获得永生。同样，那些我们称为祂仇敌的人，或仇敌的幼年子女，凡祂愿意如此重生，使他们以那藉爱德工作的信德结束此生的，他们如今（在这一切实现之前）就在那预定中成为祂的子女，并被赐予祂的圣子基督，好使他们不失丧，反而获得永生。\par

% structure: p0049 source=h2 target=section reason=relative-heading-rank; base=h2

\section{第22章——天主的真子女是基督的真门徒。}

最后，救主亲自说：\BibleQuote{你们如果存留在我内，才真是我的门徒。}（\BibleRef{若 8:31}）那么，犹大能算在他们中间吗？因为他没有存留在祂的话内。那些福音中如此描述的人能算在他们中间吗？当主命令吃祂的肉、喝祂的血时，福音记载说：\BibleQuote{这些话是耶稣在葛法翁会堂教训人时说的。祂的门徒中有许多听了，就说：\InnerQuote{这话生硬，谁能听得进去呢？}耶稣自知门徒对他这话窃窃私议，便向他们说：\InnerQuote{这话使你们起反感吗？那么，如果你们看见人子升到祂原先所在的地方，又将怎样呢？使人生活的是神，肉一无所用；我给你们所讲论的话，就是神，就是生命。但你们中间有些人却不相信。}原来耶稣从起头就知道哪些人不信，和谁要出卖祂。所以祂又说：\InnerQuote{为此，我对你们说过：除非蒙父恩赐，谁也不能到我这里来。}从此，祂的门徒中有许多人退去了，不再和祂同行。}（\BibleRef{若 6:59-66}）难道福音的话不也称这些人为门徒吗？然而他们不是真门徒，因为他们没有存留在祂的话内，正如祂所说：\BibleQuote{你们如果存留在我内，才真是我的门徒。}（\BibleRef{若 8:31}）因此，因为他们没有恒心，不真是基督的门徒，所以他们即使表面上像是天主的子女，也被称为天主的子女，但实际上不是真正的子女。我们称那些重生后我们看到虔敬生活的人为被选者、基督的门徒、天主的子女；但如果他们恒心持守那使他们如此被称的原因，他们才是真正的。但如果他们没有恒心——也就是说，他们不持守他们所开始成为的——那么他们称为的就不是真正所是的；因为在天主眼中，他们不是，祂知道他们将要成为什么——也就是说，从善人变为恶人。\par

% structure: p0051 source=h2 target=section reason=relative-heading-rank; base=h2

\section{第23章——唯独按旨意蒙召的人是被预定的。}

为此，宗徒说：\BibleQuote{我们知道：一切事都为爱天主的人带来益处}——他知道有些人爱天主，却不能在这善道上坚持到底——立刻加上：\BibleQuote{就是那些按祂旨意蒙召的人。}（\BibleRef{罗 8:28}）因为这些人对天主的爱持续到终；而那些暂时偏离正道的，也会回头，以便持续到底他们所开始成为的善。然而，为表明什么是按祂旨意蒙召，他随即加上我已在上文引述的话：\BibleQuote{因为祂所预知的人，也预定他们与祂圣子的肖像相同，好使祂在众多弟兄中作长子。祂所预定的人，也召叫了他们；祂所召叫的人，也使他们成义；祂所成义的人，也使他们得光荣。}（\BibleRef{罗 8:29-30}）这一切都已成就：祂预知、预定、召叫、成义；因为所有人都已被预知和预定，且许多人已蒙召和成义；但祂最后说的\BibleQuote{使他们得光荣}（如果这里的光荣是指同一宗徒所说的：\BibleQuote{当基督、你们的生命显现时，那时你们也要与祂一同在光荣中出现。}（\BibleRef{哥 3:4}））尚未实现。虽然这两件事——即祂召叫和祂成义——尚未在所有被说的人身上完成（因为直到世界终了，仍有许多人等待被召叫和成义），然而祂却用了过去时态来论及未来的事，仿佛天主从永恒中已安排它们发生。为此，依撒意亚先知论到祂说：\BibleQuote{那创造了未来之事的那一位。}（\BibleRef{依 45:11}）因此，凡是在天主最明智的安排中被预知、预定、召叫、成义、得光荣的人——我不仅是说尚未重生，甚至尚未出生的人——已经是天主的子女，且绝对不能丧亡。这些人真正来到基督面前，因为他们是以祂亲自说的方式来的：\BibleQuote{凡父赐给我的人，必到我这里来；到我这里来的，我决不抛弃他。}（\BibleRef{若 6:37}）稍后祂又说：\BibleQuote{派遣我来的父的旨意是：凡祂赐给我的，叫我一个也不失掉。}（\BibleRef{若 6:39}）因此，恒心行善直到终也是由祂赐下的；因为只赐给那些不会丧亡的人，因为不恒心的人必将丧亡。\par

% structure: p0053 source=h2 target=section reason=relative-heading-rank; base=h2

\section{第二十四章——即便是蒙选者的罪恶，天主也使之转为他们的益处。}

对于爱祂的人，天主使一切协助他们获得益处；真是\Emph{一切}事，以致即便他们中有人迷途、偏离正路，甚至这件事本身，祂也使之成为他们的益处，好使他们更谦卑、更有智慧地返回。因为他们知道，即使在正道上，也应当战战兢兢地喜乐；不可将持守的信心归于自己，仿佛依靠自己的力量；不可在顺境中说：\Quote{我们永远不会被动摇。}因此，经上对他们说：\Quote{当存畏惧事奉主，当战战兢兢地向祂欢乐，免得主发怒，你们就迷失了正路。}因为祂不是说：\Quote{你们不可进入正路；}而是说：\Quote{免得你们从正路上灭亡。}这难道不是表明，那些已经在正路上行走的人，被提醒要存畏惧事奉天主；也就是说，\Quote{不可心高气傲，反要战兢}？\BibleRef{罗 11:20}这意思是，他们不可骄傲，而应谦卑。此外，祂在另一处也说：\Quote{不可心高气傲，反要俯就卑微的人。}\BibleRef{罗 12:16}让他们在主内喜乐，却要战战兢兢；不可夸耀任何事物，因为没有什么属于我们自己，好叫夸耀的，因主而夸耀，免得他们从已经开始行走的正道上灭亡，因为他们将自己仍在正道上的事实归于自己。宗徒也使用这些话，说：\Quote{你们要怀着恐惧战栗，努力成就你们的救恩。}\BibleRef{斐 2:12}他又说明为什么要恐惧战栗：\Quote{因为是天主在你们内工作，使你们愿意，并使你们力行，为成就祂的善意。}\BibleRef{斐 2:13}因为那在顺境中说\Quote{我永远不会被动摇。}的人，并没有这种恐惧战栗。但因他是恩许之子，不是丧亡之子，他在天主暂时离弃他时，体验了自己究竟是谁：\Quote{主啊，}他说，\Quote{在祢的恩宠中，祢赐给了我力量，使我尊荣；祢转面不顾我，我就惊慌失措。}看！他受了多大的教训，也因此而更加谦卑地行走在他的路上，终于看见并承认，是天主凭自己的旨意，赐给他力量以维持他的尊荣。而他曾将这一切归功于自己，并以为是出于自己，好像天主所赐的丰盛是从己出，而不是从赐予者而来，因此他曾说：\Quote{我永远不会被动摇！}因此他惊慌失措，以致认识了自己，并因谦卑之心，不仅学会了关于永生的事，更学会了关于此生中虔诚的品行和坚忍，以及如何在其中保持希望。这也可能是宗徒\Person{伯多禄}的话，因为他也在顺境中说：\Quote{我要为祢舍掉我的性命。}\BibleRef{若 13:37}因着热忱，他将后来由主所赐予他的事归给自己。但主转面不顾他，他就惊慌，以至因怕为祂而死，三次否认了祂。但主再次转面向他，用他的眼泪洗去了他的罪。\Quote{祂转过身来看他}\BibleRef{路 22:61}不是别的，正是主将先前暂时转离他的面容重新赐给他吗？因此他惊慌，但因为他学会了不要信赖自己，甚至这件事也为他大有益处，这是出于那位使万事协助那些爱祂的人获得益处的大能；因为他已按旨意蒙召，以致没有人能将他从基督手中夺去，他已被赐给了基督。\par

% structure: p0055 source=h2 target=section reason=relative-heading-rank; base=h2

\section{第二十五章——因此当使用谴责。}

因此，谁也不可说，当一个人偏离正路时，不该受谴责，而只该为主祈求他回头并持守。任何明理而有信德的人不可这样说。因为这样的人若按旨意蒙召，毫无疑问，甚至在他受谴责这件事上，天主也在协助他获得益处。但由于谴责的人不知道他是否这样蒙召，就让他以爱德做他知道当做的事；因为他知道这样的人应当受谴责。天主将显示仁慈或审判：仁慈，确实，若受谴责者因恩宠的赋予而\Quote{被区别出来}，脱离了沉沦的群众，不属于那为毁灭而完成的义怒之器，而属于天主为光荣所预备的仁慈之器\BibleRef{罗 9:22}；但审判，若他属于前者而被定罪，而不蒙预定归于后者。\par

% structure: p0057 source=h2 target=section reason=relative-heading-rank; base=h2

\section{第二十六章（第十章）——亚当是否领受了恒心之恩。}

这里又生出另一个问题，不可轻率放过，而应藉着主的助佑加以探讨和解决，因为我们和我们的言论都在祂手中。\BibleRef{智 7:16} 因为有人问我，关于这个天主的恩赐——即至终坚持于善，我如何看待原祖本身，他确是无瑕地被造正直。我并非说：\Quote{若他没有坚持，他如何是无瑕的，既然他缺少如此必要的天主的恩赐？}因为对这个提问，答案是容易的：他没有坚持，因为他没有在那无罪的善中坚持；从他跌倒之时起，他便开始有罪；若他开始有罪，那么在他开始之前他当然是无罪的。因为无罪是一回事，而安处于无罪之善中则是另一回事。正是由于他并非从未无罪，而是没有继续无罪，这无疑表明他曾无罪，因为他被责备没有继续在那善中。但更应费心探讨和解答的是，我们如何回答那些说这样的话的人：\Quote{若他在那被造无罪的正直中拥有坚持，那么他无疑坚持了；若他坚持，他当然没有犯罪，也没有离弃他的正直。但他犯了罪，且离弃了善，这是真理所宣告的。因此他在那善中没有坚持；若他没有坚持，他当然没有领受它。因为他怎能既领受了坚持，又没有坚持呢？再者，若他没有坚持是因为他没有领受坚持，那么他因没有坚持而犯了什么罪呢？因为他不能说是为这原因而没有领受：他没有藉恩宠的施予从毁灭的群体中被分别出来，因为在那个从他而出的败坏根源的人犯罪之前，毁灭的群体尚未在人类中存在。}\par

% structure: p0059 source=h2 target=section reason=relative-heading-rank; base=h2

\section{第27章 答复。}

因此，我们最有益地承认我们所最正确相信的：万有的天主和主，祂以祂的权能创造了所有美好的事物，并预知邪恶将从善中产生，且知道祂那全能的美善更愿意从恶中行善，而非不允许任何恶存在，祂如此安排了天使和人的生命，在其中首先显明了他们的自由意志能做什么，然后显明了祂恩宠的仁慈和祂公义的审判能做什么。最后，某些天使，其中为首的被称为魔鬼，藉自由意志背离了主天主。然而，虽然他们逃离了祂的良善（在其中他们曾是幸福的），却不能逃离祂的审判，藉此他们变得极其不幸。但其他天使，藉同样的自由意志坚守于真理，并堪当知道那最确实的真理：他们永不会堕落。因为若我们从圣经得以知道没有任何圣天使会再堕落，那么他们自己藉更崇高地向他们启示的真理岂不更知道这事？因为应许给我们的是无尽的幸福生命，并等同于天使，\BibleRef{玛 22:30} 由此应许我们确信，当我们在审判之后到达那生命时，我们将不会从其中堕落；但若天使对自己不知道这真理，我们就不会等同于他们，反比他们更有福。但真理应许我们等同于他们。因此，他们已藉看见知道这事，而我们藉信德知道，即如今再没有圣天使的堕落。但魔鬼及其天使，虽然他们在堕落之前曾是幸福的，且不知道他们会堕入悲惨——仍有某种东西可以加于他们的幸福，若他们藉自由意志坚守于真理，直到他们领受那作为坚持酬报的最高幸福之圆满；即藉着圣神所赐的天主之爱的极大丰盈，他们绝对不能再堕落，并且他们对自己完全确知此事。他们没有这份圆满的幸福；但既然他们不知道自己未来的悲惨，他们享受一种较小但仍无任何缺陷的幸福。因为他们若知道自己未来的堕落和永罚，他们当然不能幸福；因为对如此巨大恶事的恐惧甚至那时就会使他们不幸。\par

% structure: p0061 source=h2 target=section reason=relative-heading-rank; base=h2

\section{第28章 原祖本人也可能凭自由意志坚持。}

如此，祂也以\OriginalTerm{自由意志}{free will}造了人；人虽不知自己未来的堕落，却因此幸福，因他以为不死亡、不陷于不幸都在自己的能力之内。若他愿意凭自己的自由意志继续处于这种正直与无罪状态，确实，完全无需经历死亡与不幸，他就会因这种坚持的功劳而获得圆满的祝福，就如圣天使们所蒙受的祝福一样；也就是说，再也不能堕落，并且绝对确知这点。因为即使他自己在乐园中也不能蒙福——不，他若因预知自己的堕落而恐惧如此灾祸，从而痛苦，他就不会在那里，因为那里不适于他受苦。但因他凭自由意志离弃了天主，他便体验到天主公义的审判：他与他的整个族类——当时全在他内，与他一同犯了罪——应被定罪。因为凡从这族类中蒙受天主恩宠而得救的，确实都是从他们已被束缚的定罪中被拯救出来的。因此，即使没有一个人得救，无人能公正地指责天主的审判。因此，与那些丧亡者相比，得救者\Emph{稀少}，但就其绝对数目而言，却是\Emph{众多}；这是由恩宠成就的，是由白白恩赐成就的：必须感恩，因为这是成就的，好使无人因自己的功绩而自夸，而是使\Quote{各种借口都被塞住}\BibleRef{罗 3:19}，且\Quote{夸耀的，应因主而夸耀}\BibleRef{耶 9:24}。\par

% structure: p0063 source=h2 target=section reason=relative-heading-rank; base=h2

\section{第二十九章［XI］——堕落前与堕落后所赐恩宠的区别}

那么，难道亚当没有领受天主的恩宠吗？是的，确实，他丰沛地领受了，但种类不同。他被安置在从其创造者的美善所领受的诸多恩惠之中；因为他并未凭自己的功德获得这些恩惠；在这些恩惠中，他完全不受任何邪恶。但此世的圣人们——他们蒙受了这解救的恩宠——却处于诸多邪恶之中，他们向天主呼求说：\Quote{救我们免于凶恶}\BibleRef{玛 6:13}。他在那些恩惠中不需要基督的死亡；这些人，那羔羊的血却洗除了他们承继而来的和自身的罪债。他不需要那援助——当他们说\Quote{我发现在我的肢体中另有一条法律，与我心神的法律交战，把我掳去，叫我隶属于那在我肢体中罪恶的法律。我这个人真不幸呀！谁能救我脱离这取死的身体呢？天主的恩宠，藉着我们的主耶稣基督}\BibleRef{罗 7:23}时，他们所恳求的援助。因为在他们内，肉身与心神相争，心神与肉身相争；当他们在这场争战中劳苦并遇险时，他们祈求藉着基督的恩宠赐予他们战斗与得胜的力量。而他，却从未在类似的自我内在冲突中受试探、受困扰，在那福乐的状态中安享与自身的平安。\par

% structure: p0065 source=h2 target=section reason=relative-heading-rank; base=h2

\section{第三十章——圣言的降生成人}

因此，虽然这些人现在不需要更喜乐的恩宠，他们却需要更强大的恩宠；有什么恩宠比天主的独生子——与圣父同等、同永恒，为他们而成为人，且自身毫无任何原罪或本罪，却被罪人钉在十字架上——更强大呢？虽然祂在第三日复活，永不再死，但祂却为人死，并赐生命给死者，好使被祂的血救赎、领受了如此伟大保证的他们能说：\Quote{若是天主偕同我们，谁能反对我们呢？祂没有怜惜自己的儿子，反而为我们众人把祂交出，这样既没有把祂连祂一切赐给我们吗？}\BibleRef{罗 8:31}因此，天主取了我们的本性——即基督的人的灵魂与肉身——藉着一种极其奇妙或奇妙至极的承担；如此一来，祂没有任何自己义德的先行功德，却从起初——祂开始成为人之时——就如此是天主子，以致祂与那无始的圣言成为同一\OriginalTerm{位格}{person}。因为没有人对此事及信仰如此盲目，竟敢说：虽然由圣神及童贞玛利亚所生的人子，却凭自己的自由意志，藉着义德的生活和行善而毫无罪过地配成为天主子；这是相反福音的，福音说：\Quote{圣言成了血肉}\BibleRef{若 1:14}。因为除了在童贞女的子宫中——那是基督之人的开始——祂在那里成了血肉，又在哪里成了血肉呢？此外，当童贞女询问天使所告诉她的这事如何成就时，天使回答说：\Quote{圣神要临于你，至高者的能力要荫庇你，因此，那要诞生的圣者，将称为天主子}\BibleRef{路 1:35}。\Quote{因此}，他说；不是因为善行——因为在尚未出生的婴儿身上，当然没有善行；而是\Quote{因此}，因为\Quote{圣神要临于你，至高者的能力要荫庇你，那要诞生的圣者，将称为天主子}。那诞生绝对是无偿的，在位格的合一中将人与天主、血肉与圣言结合！善行随着那诞生而来；善行并未赚得那诞生。因为绝无必要担心那被天主圣言以如此不可言喻方式取为位格合一的人性会藉自由意志的选择而犯罪，因为那取用本身正是如此，以致如此被天主取用的人性不会接纳任何恶意的活动。天主藉此中保表明，祂使那些被祂的血从邪恶中救赎的人永远为善；而祂以如此方式取用了祂，使祂永不成为恶的，并且既非从恶中造成，就永远为善。\par

% structure: p0067 source=h2 target=section reason=relative-heading-rank; base=h2

\section{第三十一章——第一个人已领受了为恒心所必需的恩宠，但其运用却留在了他的自由选择中}

第一个人没有那种使他永远不会愿意作恶的恩宠；但他确实有那样的恩宠，如果坚持在其中，他就永远不会作恶，而且没有这恩宠，他凭自由意志不能为善，然而他凭自由意志却能离弃这恩宠。因此，天主甚至不愿意他缺少自己的恩宠，天主将这恩宠留在他自由意志中；因为自由意志足以作恶，但要为善则不足，除非受全能之善的援助。如果那个人没有离弃他自由意志的这种援助，他就会永远为善；但他离弃了它，也就被离弃了。因为这样的援助其本性是：他可以在愿意的时候离弃它，也可以在愿意的时候继续在其中；但却不能使他必然愿意。这首先是赐给第一个亚当的恩宠；但第二个亚当的恩宠比这更强大。因为前者是使人若愿意就能拥有正义；而后者则能做得更多，因为它甚至使人愿意，并且愿意得如此强烈，爱得如此炽热，以至于藉着圣神的意愿战胜那与之相争的肉身的意愿。确实，那恩宠也不算小，它甚至显示了自由意志的能力，因为人得到如此帮助，以致没有这帮助就不能坚持为善，但可以离弃这帮助，如果他愿意。但这后来的恩宠更加伟大，因为人藉着它恢复失去的自由还远远不够；最终，没有它，人既不能领会善，也不能在善中坚持（如果他愿意），除非他也被\Emph{促使愿意}。\par

% structure: p0069 source=h2 target=section reason=relative-heading-rank; base=h2

\section{第三十二章——在创造中赐予亚当的恩宠}

那时，天主赐给了人一个好的意志，因为祂造人时使人正直，就赐予了那意志。祂赐予了帮助，没有这帮助，人若愿意就不能继续在其中；但祂将愿意与否留在了人的自由意志中。因此，人若愿意就能继续，因为不缺那使他能继续的援助，没有这援助他就不能坚忍地持守他所愿意的善。但他不愿意继续，这完全是人的过错，因为他若愿意继续，本应成为他的功德；正如圣天使那样，其他天使因自由意志而堕落，他们却因同一自由意志而站立，并因这继续而配得应得的赏报——即如此丰满的福乐，使他们能极其确定地永远安于其中。然而，如果这种帮助在起初的天使或人身上是缺乏的，因为他们的本性并非被造得在缺少神圣帮助时还能（若愿意）坚持，那么他们必定不会因自己的过错而堕落，因为那使他们不能继续的帮助是缺失的。但现在，对于那些缺少这种援助的人，这是罪的惩罚；但对于那些得到帮助的人，这是恩宠的赐予，而非债务；而且通过我们的主耶稣基督，它更丰盛地赐给天主愿意赐予的人，以致我们不仅有那帮助（即没有它即使愿意也不能继续），而且还有如此伟大和如此大的帮助，以至于我们\Emph{愿意}。因为藉着天主的这恩宠，在我们领受善和坚忍持守善时，不仅使我们能行我们所愿的，甚至使我们愿意行我们所能的。但第一个人的情况并非如此；因为他有前者，却没有后者。因为他不需要恩宠来领受善，因为他尚未失去它；但他需要恩宠的帮助来继续在其中，没有这帮助他根本无法做到；他获得了能力（如果他愿意），但他没有意愿去做他能做的事；如果他有了意愿，他就会坚持。因为他若能愿意就能坚持；但不愿意却是自由意志的结果，那自由意志当时如此自由，以至于他能同时愿意善和恶。因为还有什么比那不能服事罪的自由意志更自由呢？这也应当临到人身上，就如它被造成圣天使的赏报一样。但现在，由于罪而失去了善的功德，在被拯救的人身上，那本应是功德赏报的，成了恩宠的礼物。\par

% structure: p0071 source=h2 target=section reason=relative-heading-rank; base=h2

\section{第三十三章 [XII.]——什么是能够不犯罪、能够不死、能够不离弃善的能力与不能犯罪、不能死、不能离弃善的能力之间的区别？}

因此，我们必须仔细并注意地思考这几对之间的区别——能够不犯罪与不能犯罪；能够不死与不能死；能够不离弃善与不能离弃善。因为第一个人能够不犯罪，能够不死，能够不离弃善。我们能说那拥有如此自由意志的人不能犯罪吗？或者能对那被警告\Quote{你若犯罪，就必死}的人说不能死吗？或者能说他不能离弃善，而他却因犯罪而离弃善，并因此死亡？因此，意志的第一自由是\Emph{能够不犯罪}，最后的自由将伟大得多，是\Emph{不能犯罪}；第一个不朽是\Emph{能够不死}，最后的将伟大得多，是\Emph{不能死}；第一种坚忍的能力是\Emph{能够不离弃善}，最后的将是\Emph{不能离弃善}的福乐。但因为最后的祝福更优越、更好，那么最初的那些祝福就因此要么根本不是祝福，要么是微不足道的祝福吗？\par

% structure: p0073 source=h2 target=section reason=relative-heading-rank; base=h2

\section{第三十四章——没有它事情就不会发生的援助，和拥有它事情就会发生的援助}

此外，这些援助本身应当加以区分。没有援助事情就不会发生，这是一回事；有了援助事情才得以发生，则是另一回事。因为没有食物我们就无法存活；然而，即使食物在手，若一个人愿意去死，食物也不能使他存活。因此，食物的援助是那种没有它我们就无法存活的援助，而不是那种使我们存活的援助。但是，当一个人获得了他所没有的真福时，他立刻就成为了有福的。因为这种援助不仅是那种没有它事情就不会发生的援助，而且也是那种有了它事情就得以发生、并为此目的而赐予的援助。因此，这是一种既使之发生、又不可或缺的援助；因为一方面，如果真福赐予一个人，他立刻就成为了有福的；另一方面，如果从未赐予，他也永远不会成为有福的。但食物并不必然使人存活，尽管没有食物他就无法存活。因此，对于最初的那个人，他在被造为正直的善之中，已经接受了不犯罪的能力、不死的能力、不抛弃那善本身的能力，他被赐予了恒心的援助——不是那种使之恒心的援助，而是那种没有它他就无法凭自由意志恒心的援助。但现在，对于被天主的恩宠预定进入天主之国的圣人们，所赐予的恒心援助并非像前者那样，而是将恒心本身赐予他们；不仅使他们没有此恩赐就无法恒心，而且使他们藉此恩赐不得不恒心。因为祂不仅说：\BibleQuote{离了我，你们什么也不能做，}\BibleRef{若 15:5}而且祂还说：\BibleQuote{不是你们拣选了我，而是我拣选了你们，并委派你们去结果实，叫你们的果实常存。}\BibleRef{若 15:16}这些话表明，祂不仅赐予了他们义德，也赐予了他们其中的恒心。因为当基督如此委派他们去结果实，并使他们的果实常存时，谁敢说它不会常存？谁敢说它或许不会常存？\BibleQuote{因为天主的恩赐和恩召是不会撤回的。}\BibleRef{罗 11:29}但恩召是针对那些按照祂旨意蒙召的人。因此，当基督为这些人代祷，使他们的信仰不致丧失时，无疑它将不会丧失直到最后。这样，它将恒心到底，而此生的终结也不会发现它有任何减退。\par

% structure: p0075 source=h2 target=section reason=relative-heading-rank; base=h2

\section{第三十五章——现在圣人们比在亚当内有更大的自由。}

面对如此多且大的诱惑——这些在乐园里并不存在——确实需要更大的自由，它由恒心的恩赐所坚固和确认，以便战胜这个世界及其一切爱慕、恐惧和谬误：圣人们的殉道证明了这一点。最终，他（亚当）不仅无人使之畏惧，而且反之，尽管有天主权威的恐惧，他使用自由意志，却没有在这种幸福状态和如此容易地不犯罪中站立得住。但这些（圣人们），我说，不是在世界的恐惧之下，而是不顾世界的怒火（惟恐他们站立），却坚定地站立在信仰中；而他（亚当）能看到将要舍弃的现世美物，他们（圣人们）却不能看到将要领受的未来美物。这是从何而来？若非出于祂的恩赐，他们从祂获得了怜悯以致成为忠信者；从祂领受了不是惧怕的灵——否则他们就会向迫害者屈服——而是刚毅、爱德和节制的灵，藉此他们能胜过一切恐吓、诱惑和折磨。因此，他（亚当）虽无任何罪，却被赐予了受造时的自由意志，而他使这意志服侍了罪。但这些（圣人们）的意志虽曾是罪的奴仆，却被那位说过\BibleQuote{若是子使你们自由，你们就真自由了。}\BibleRef{若 8:36}的所释放。藉此恩宠，他们获得如此大的自由，以至于尽管他们在此生中一直与情欲的罪恶争战，且有一些罪恶不知不觉地潜入，为此他们每日念诵\BibleQuote{宽免我们的罪债，}\BibleRef{玛 6:12}但他们不再顺从那至于死的罪——关于这罪，若望宗徒说：\BibleQuote{有至于死的罪，我不说要为这罪祈求。}\BibleRef{若一 5:16}关于这罪（既然未明说），可以有许多不同的理解。然而，我说，这罪就是至死背弃那藉爱德工作的信仰。那些不在起初的状态（如亚当一样自由）、而是藉第二位亚当蒙受天主的恩宠被释放的人，不再服侍这罪；藉着那释放，他们拥有那使他们能服侍天主的自由意志，而不是那使他们可能被魔鬼俘虏的自由意志。他们从罪中得释放，成了义德的奴仆，\BibleRef{罗 6:18}在这义德中他们将坚持到底，这是藉着从祂而来的恒心恩赐——祂预知他们、预定他们、按照祂旨意召叫他们、使他们成义、并光荣了他们，因为祂已经形成了祂所应许关于他们的未来之事。当祂应许时，\BibleQuote{亚巴郎相信了祂，这就算为他的义德。}因为\BibleQuote{他完全相信，将光荣归于天主，}正如经上所载：\BibleQuote{祂所应许的，也必能成就。}\par

% structure: p0077 source=h2 target=section reason=relative-heading-rank; base=h2

\section{第三十六章——天主不仅预知人将成为善的，而且祂亲自使他们成为善的。}

因此，正是祂自己使那些人成为好人，以行善工。因为祂向亚巴郎应许他们，并非由于祂预知他们会凭自己成为好人。若是这样，祂所应许的便不是属于祂，而是属于他们了。但亚巴郎并非如此相信，而是\BibleQuote{他信德不弱，将光荣归于天主；}并且\BibleQuote{满心相信，祂所应许的，祂也必能成就。}（\BibleRef{罗 4:19}）祂并不是说：\Quote{祂预知的，祂能应许；}也不是说：\Quote{祂预言的，祂能显明；}也不是说：\Quote{祂应许的，祂能预知；}而是说：\Quote{祂应许的，祂也必能成就。}因此，是祂使他们持守于善，祂使他们成为好人。但那些堕落和丧亡的人从不属于被预定者。虽然如此，当宗徒论及所有重生且虔诚生活的人时说：\BibleQuote{你是谁，竟论断别人的仆人？他或站住或跌倒，自有他的主人；}但他立刻顾念到被预定者，说：\BibleQuote{而且他必要站住；}并且为使众人不归功于自己，他说：\BibleQuote{因为天主能使他站住。}（\BibleRef{罗 14:4}）因此，正是祂自己赐予坚忍，祂能扶持那些站住的人，使他们以最大的坚忍站稳；或恢复那些跌倒的人，因为\Quote{主扶持那被压倒的人。}\par

% structure: p0079 source=h2 target=section reason=relative-heading-rank; base=h2

\section{第三十七章。——坚固的意志被托付了坚忍或不坚忍的能力。}

因此，正如第一个人没有领受天主的这个恩赐——即持守于善的坚忍，而是让坚忍或不坚忍留在他自己的选择中，他的意志曾具有这样的力量——因为它被造时毫无罪过，也没有任何私欲反抗它——以致坚忍的选择能够相称地被托付给如此良善且易于善生的意志。但天主同时预知他会行不义；预知了，却没有强迫他行不义；同时祂也知道祂自己会以正义对待他。但现在，由于那伟大的自由因罪的功过而丧失，我们的软弱仍须靠更大的恩赐来扶助。因为天主乐意最有效地扑灭人类骄傲的妄念，\BibleQuote{使凡有血肉的，在天主面前一个也不能自夸}（\BibleRef{格前 1:29}）——即\Quote{无人}。但血肉在天主面前为何不能自夸呢？岂不是因它的功德吗？它本来可以有功德，却丧失了；并且恰恰是通过它本可能拥有功德的那方式丧失的，即通过它的自由意志；为此，那些将要得救的人，除了救主的恩宠之外，一无所有。因此，没有血肉在天主面前自夸。因为不义的人不自夸，因为他们没有可夸的；义人也不自夸，因为他们的夸耀来自祂，他们不归功于自己，而归功于祂，他们对祂说：\Quote{我的荣耀，使我抬起头来的。}因此，经上所记的适用于每一个人：\Quote{使凡有血肉的，在天主面前一个也不能自夸。}然而，义人适用这段经文：\BibleQuote{夸耀的，应当因主而夸耀。}（\BibleRef{格前 1:31}）因为宗徒最清楚地表明这一点，当他说\Quote{使凡有血肉的，在天主面前一个也不能自夸}之后，以免圣人们以为他们毫无可夸之处，他立刻加上：\BibleQuote{但你们得以在基督耶稣内，是由于祂，祂由天主成了我们的智慧、正义、圣化与救赎：好使按经上所载的：夸耀的，应当因主而夸耀。}（\BibleRef{格前 1:30}）因此，在这苦难的居所，人生在世就是考验，\BibleQuote{能力在软弱中才得以成全}（\BibleRef{格后 12:9}）。什么能力？岂不是\Quote{那夸耀的，应当因主而夸耀}？\par

% structure: p0081 source=h2 target=section reason=relative-heading-rank; base=h2

\section{第三十八章。——如今赐予圣徒的坚忍之恩的本质是什么。}

于是，天主愿意祂的圣人们——即使是关于\OriginalTerm{恒心}{perseverance}本身——也不要在自己的能力中夸耀，而要在祂内夸耀；祂不仅赐给他们如同赐给第一个人的援助，没有这援助他们即便愿意也无法恒心，祂也在他们内引发意志；因为，除非他们既能又愿意，否则他们不会恒心，因此恒心的能力和意愿都应藉着天主恩宠的慷慨赐予他们。因为藉着圣神，他们的意志被如此点燃，以致他们因此能够，因为他们如此愿意；他们因此如此愿意，因为天主在他们内运作，使之愿意。因为如果在此生如此众多的软弱中（然而，在此软弱中，为了抑制骄傲，力量理当被成全），他们自己的意志被留给他们自己，以便他们如果愿意，就可以继续在天主的帮助中（没有这帮助他们不能恒心），而天主不在他们内运作使之愿意，那么在如此众多且如此重大的软弱中，他们的意志本身就会屈服，他们就不能恒心，理由是：由于软弱而失败，他们就不再愿意，或者由于意志的软弱，他们不会如此愿意以致能够。因此，援助被带给人类意志的软弱，好使它被天主的恩宠不可改变且不可战胜地影响；这样，尽管软弱，它仍然可以不失败，也不被任何逆境所克服。于是，人的意志（软弱且无能，在尚小的善中）因天主的力量得以恒心；而第一个人的意志（强壮且健康，拥有\OriginalTerm{自由选择}{free choice}的能力）却在更大的善中没有恒心；因为尽管天主的帮助并不缺乏（没有这帮助，如果他愿意，他就不能恒心），但这帮助并非那种天主在其中运作使人愿意的帮助。的确，祂容许最强者做祂所愿意的事；而对于弱者，祂保留着藉祂自己的恩赐，使他们不可战胜地愿意善事，并不可战胜地拒绝放弃此善。因此，当基督说：\BibleQuote{我曾为你祈求，使你的信德不致丧失} \BibleRef{路 22:32}，我们可以理解为这是对那建立在磐石上的人说的。因此，天主的人，不仅因为他得到了怜悯而成为忠信的，而且因为信德本身不会丧失，如果他要夸耀，他必须因主而夸耀。\par

% structure: p0083 source=h2 target=section reason=relative-heading-rank; base=h2

\section{第三十九章 —— \OriginalTerm{被预定者}{Predestinated}的数目是确定且限定的}

我这样说论及那些被预定进入天主国度的人，他们的数目如此确定，以致不能增添也不能减少；不是论及那些当祂宣布和发言时、成为无数的人。因为他们可被称为蒙召却未被拣选，因为他们不是按旨意蒙召的。但被选者的数目是确定的，既不能增加也不能减少——尽管这事由\Person{洗者若翰}所预示，当他说：\BibleQuote{那么，你们就结与悔改相称的果实吧！你们自己不要思念说：我们有亚巴郎为父。我给你们说：天主能从这些石头给亚巴郎兴起子孙来。} \BibleRef{玛 3:8}，以表明他们若不结果实，就要这样被砍掉，以致那应许给\Person{亚巴郎}的数目不会缺少——却在\BibleBook{}{默示录}中更清楚地宣告：\BibleQuote{你应坚持你所有的，免得别人拿去你的冠冕。} \BibleRef{默 3:11}。因为如果\Emph{别人}不会领受，除非\Emph{一人}失去了，那么数目就是固定的。\par

% structure: p0085 source=h2 target=section reason=relative-heading-rank; base=h2

\section{第四十章 —— 没有人对自己的\OriginalTerm{预定}{Predestination}和\OriginalTerm{救恩}{Salvation}是确定和安全的}

而且，这样的言论对将要恒心的圣人们说，好像他们是否能恒心被视为不确定，这是一个理由，使他们不应不然听这些事，因为对他们有益的是：\BibleQuote{不可心高气傲，而要恐惧战兢。} \BibleRef{罗 11:20} 因为信友中的众人，当他还活在这必死的状态中时，谁能妄自假定自己是在被预定者的数目中呢？因为在此状态中，此事必须保持隐藏；因为在此我们要如此警惕骄傲，以致连一位如此伟大的宗徒也被撒殚的使者击打，免得他高举自己。 \BibleRef{格后 12:7} 因此，对宗徒们说：\BibleQuote{你们若住在我内} \BibleRef{若 15:7}，而说这话的那位确知他们必会住下；又藉先知说：\BibleQuote{如果你们愿意，且听从我} \BibleRef{依 1:19}，尽管祂也知道祂要在谁内运作使之愿意。许多类似的话也说了。因为由于这隐藏之事的益处，免得有人被高举，反而所有即使跑得好的人应当恐惧，因为不知道谁能达到——由于这隐藏之事的益处，必须相信，有些丧亡之子，没有领受恒心至终的恩赐，开始活在以爱德行事的信德中，并忠信且正义地生活一段时间，之后便跌倒，而且在发生这事之前没有被从这个生命取去。如果这事没有发生在任何人身上，人们就会拥有那种非常健康的恐惧，藉此使妄自尊大的罪被压制，但仅限于直到他们获得基督的恩宠，藉以度虔敬的生活，而后他们将来就会确信自己绝不会离弃祂。而且如此妄自尊大在这试炼的状态中是不合宜的，这里有如此巨大的软弱，以致安全感可能滋生骄傲。最后，这也将成事实；但将在那时，在人身上也如同已在天使身上一样，那时不能有任何骄傲。因此，圣人的数目，由天主的恩宠预定进入天主的国度，并获赐恒心至终的恩赐，将完全地被引导到那里，并最终将在无终的圆满中保存，最为有福，他们救主的慈爱仍紧随着他们，无论是在他们的皈依中，在他们的战斗中，还是在他们的冠冕中！\par

% structure: p0087 source=h2 target=section reason=relative-heading-rank; base=h2

\section{第四十一章 —— 即使在\OriginalTerm{审判}{Judgment}中，\OriginalTerm{天主的慈悲}{God's Mercy}对我们仍是必需的}

因为圣经证明，天主的仁慈在此时对他们也是必要的，当圣人对自己的灵魂论及它的主天主说：\Quote{祂以仁慈和怜悯给你加冕。}宗徒雅各伯也说：\BibleQuote{没有施行仁慈的，将受无仁慈的审判；}\BibleRef{雅 2:13}他在这里指出，即使在义人受冠冕、不义者被定罪的那个审判中，有些人将受仁慈的审判，另一些人将受无仁慈的审判。为此，玛加伯的母亲也对她的儿子说：\BibleQuote{好使我在那仁慈中，与你的弟兄们一同接纳你。}\BibleRef{加下 7:29}\BibleQuote{因为当正义的君王，正如经上所载，坐在宝座上时，没有恶事能抵挡他。谁能夸口说他有洁净的心？谁能夸口说他纯净无罪？}\BibleRef{箴 20:8}因此，天主的仁慈甚至在那时也是必要的，藉着这仁慈，\Quote{上主不归咎于罪的人是有福的。}但在那时，仁慈本身也要按善工的功过，在公义的审判中被分派。因为当说\Quote{没有施行仁慈的，将受无仁慈的审判}时，这清楚表明，在那些有仁慈善工的人身上，审判将带着仁慈执行；因此，那仁慈本身也要回报给善工的功过。但现在却不是这样；当时不但没有善工，反而有许多恶行在先，祂的仁慈先发制人，使人得脱诸恶——既脱离他所行的恶，也脱离他若不受天主的恩宠约束就会犯的恶，还脱离他若不从黑暗的权势中被拔出、迁到祂爱子的国里就会永远承受的恶。\BibleRef{哥 1:13}然而，就连那永生本身——它确是由于善工而给予的——也被如此伟大的宗徒称为天主的恩宠，虽然恩宠不是因功行而回报的，而是白白赐予的，所以必须毫不犹豫地承认，永生之所以被称为恩宠，是因为它回报给恩宠所赐予人的那些功绩。因为福音中所读的那句话被正确理解：\BibleQuote{恩宠上加恩宠，}\BibleRef{若 1:16}——就是说，为那些由恩宠所赐的功绩。\par

% structure: p0089 source=h2 target=section reason=relative-heading-rank; base=h2

\section{第四十二章——被弃绝的人要因另一种功过受罚}

但那些不属于这预定数目的人——无论他们还没有任何意志的自由选择，还是拥有真正自由的意志选择（因为被恩宠本身所释放）——天主的恩宠将他们带入祂的国——那么，那些不属于那最确实、最蒙福数目的人，会按他们所应得的公义受审判。因为他们要么身陷由原祖遗传而继承的罪中，带着那未藉重生除去的继承债务离开此世，要么藉其自由意志添加了其他罪过；我说，他们的意志是\Emph{自由的}，但不是\Emph{被释放的}——脱离了义，却受制于罪，被各种有害的私欲所摇荡，有些更恶，有些稍轻，但全是恶的；他们必须根据这多样性被判处不同的刑罚。或者他们领受了天主的恩宠，但却只是暂时的，没有持守到底；他们离弃，也被离弃。因为藉着他们的自由意志，既然他们没有领受恒心的恩赐，他们就被天主公义而隐藏的审判所舍弃。\par

% structure: p0091 source=h2 target=section reason=relative-heading-rank; base=h2

\section{第四十三章[十四]——责斥与恩宠并不互相排斥}

因此，让人在犯罪时接受责斥，不要从责斥本身推断出反对恩宠的结论，也不要从恩宠推断出反对责斥的结论；因为罪应受公义的惩罚，公义的责斥也属于罪，如果它被用作医治，即使病人的得救是不确定的；这样，如果受责斥的人属于预定之数，责斥对他就是有益的良药；如果他不属于那数目，责斥对他就是惩罚性的打击。因此，正是在这种不确定之下，责斥必须出于爱而施行，尽管其结果未知；并且必须为受责斥的人祈祷，使他获得痊愈。但当人们藉着责斥进入或返回义路时，谁在他们心里成就救恩呢？岂不正是那赐予增长的天主吗？无论谁栽种、浇灌，无论谁在田地或灌木上劳作——当祂愿意赐予救恩时，没有人的意志能抗拒祂？因为愿意或不愿意，在于那愿意或不愿意者的权能，但并不能阻碍天主的旨意或战胜天主的权能。因为甚至对那些做祂所不愿之事的人，祂自己也做祂所愿的事。\par

% structure: p0093 source=h2 target=section reason=relative-heading-rank; base=h2

\section{第四十四章——天主愿意所有的人都得救的方式}

至于所写的，\BibleQuote{祂愿意所有的人都得救，}\BibleRef{弟前 2:4}然而并非所有的人都得救，这可以有多种理解，我在其他著作中提到过一些；但这里我要说一点：\Quote{祂愿意所有的人都得救}是这样说的，以致可以理解为所有预定的人，因为各类人都在其中。正如对法利塞人所说的：\BibleQuote{你们捐献十分之一的各样香料，}\BibleRef{路 11:42}这句话只能理解为他们所拥有的各种香料，因为他们并没有捐献遍及全地的每一种香料。按照同样的说法，\BibleQuote{好像我也在一切事上使众人喜悦。}\BibleRef{格前 10:33}因为说这话的人难道也使迫害他的群众喜悦吗？但他喜悦的是在基督的教会里聚集的各类人，无论是已经建立其中的，还是将要被引入的。\par

% structure: p0095 source=h2 target=section reason=relative-heading-rank; base=h2

\section{第四十五章——圣经实例证明天主对人的意志拥有比人自己更大的权能}

因此，毫无疑问，人的意志不能抗拒天主的旨意，以致阻止祂成就祂所愿意的事；天主\Quote{\BibleQuote{在天上地下，凡祂所喜悦的，都已作成}}祂也\Quote{\BibleQuote{成作了那些将要来的事}}\BibleRef{依 45:11}；因为祂甚至在人自己的意志上，也随己意而行，当祂愿意时。除非（仅举数例），当天主愿意将王国赐给\Person{撒乌耳}时，\Person{以色列}人是否能够（正如确实在于他们的意志）或不服从那个人，以致他们甚至能胜过天主？然而，天主成就这事，并非不藉着人自己的意志，因为祂无疑拥有最全能的大能，能随祂所喜悦地扭转人心。因为经上这样记载：\Quote{\Person{撒慕尔}遣散了百姓，各人回自己的地方去了。\Person{撒乌耳}也回\Place{基贝亚}的家去了；有心中被主感动的人，同他一起去了。但有些无赖之辈说：\InnerQuote{这人怎能拯救我们？}就鄙视他，没有给他献礼。}有谁能说，那些心被主感动而与\Person{撒乌耳}同行的人中，会有人不与他同去？或者那些未被主感动而这样做的人中，会有人去呢？至于\Person{达味}，主预定他在更兴旺的后裔中作王，我们读到：\Quote{\BibleQuote{\Person{达味}日渐强盛，天主与他同在。}}\BibleRef{编上 11:9} 这话说后不久，又记载：\Quote{\BibleQuote{那时，三杰之首\Person{阿玛赛}被圣神充满，说：\InnerQuote{\Person{达味}，我们是属你的；\Person{叶瑟}之子，我们拥护你！愿你平安，愿你的助手也平安！因为你的天主帮助了你。}}}\BibleRef{编上 12:18} 他能抗拒天主的旨意，而不做那藉着圣神（他因此被充满）在他心中运行、使他愿意、说话和这样做的那一位的旨意吗？此外，稍后同一段圣经又说：\Quote{\BibleQuote{这一切战士都怀着平和的心来到\Place{赫贝龙}，立\Person{达味}作全\Place{以色列}的王。}}\BibleRef{编上 12:38} 当然，是他们凭自己的意志立\Person{达味}为王。谁看不见这一点？谁能否认？因为他们并非被迫或非本意地这样做，而是怀着平和的心。然而，成就这事的是那在人心中随己意运行的那一位。因此，圣经先写道：\Quote{\BibleQuote{\Person{达味}日渐强盛，全能的天主与他同在。}}这样，与他同在的全能天主，感动这些民立他为王。祂如何感动他们？难道是用身体的锁链强迫他们吗？祂在内里运行；祂握住他们的心；祂激动他们的心，并藉着祂亲自在他们心中运行的意志引领他们。那么，既然当天主愿意在地上立王时，祂对人的意志拥有比人自己更大的权能，那么，还有谁能使劝诫有益，使受劝诫之人的改正结出果实，使他能在天国中坚定不移呢？\par

% structure: p0097 source=h2 target=section reason=relative-heading-rank; base=h2

\section{第46章——劝诫须依过犯的多样性而变化。教会中无大于绝罚的惩罚。}

因此，让那些在下的弟兄受那些管辖他们的人的责罚，这些责罚出自爱，根据过犯的大小而有所不同。因为那种被称为定罪、由主教判决所施加的惩罚，在教会中没有比这更大的了，如果天主愿意，可能产生并有利于最有益的责罚。因为我们不知道明日会发生什么；也不应在今世结束前对任何人绝望；也不能反驳天主，免得祂不垂顾并赐予悔改，并接受困苦的精神和痛悔之心的祭献，并免除定罪的罪债，无论多么公正，从而祂自己不审判那被定罪的人。然而牧职的需要要求，为了可怕的传染不致蔓延到多人，病羊应与健康的羊分离；或许，正是通过这种分离，由那无所不能者医治。因为我们不知道谁属于\OriginalTerm{预定}{predestinated}的数目，我们应当这样被爱的情感所影响，以致愿意所有的人得救。因为当我们努力引导每一个人到达那个点，在那里他们可能遇到那些媒介，通过我们得以成功，达到结果，即\OriginalTerm{因信成义}{justified by faith}而与天主相和，\BibleRef{罗 5:1}——这和平，此外，宗徒宣布时曾说：\BibleQuote{因此，我们为基督做大使，好像天主借着我们劝勉，我们替基督求你们与天主和好。}\BibleRef{格后 5:20} 因为\Quote{和好}是什么，不就是与祂相和吗？为了这和平，主耶稣基督亲自对祂的门徒说：\BibleQuote{你们无论进入哪一家，先要说：愿这一家平安。那里若有和平之子，你们的平安就必临到他；不然，就必归回你们。}\BibleRef{路 10:5} 当他们传扬这和平的福音——关于这福音曾预言说：\BibleQuote{那传报平安、传报佳音者的脚是何等美丽！}\BibleRef{依 52:7}——对我们而言，确实，每一个服从并相信这福音的人，就开始成为\OriginalTerm{和平之子}{son of peace}，并且因信成义，开始与天主有和平；但按照天主的预定，他早已是和平之子。因为并没有说：你们的平安临到谁，谁就必成为和平之子；但基督说：\BibleQuote{那里若有和平之子，你们的平安就必临到那家。} 所以，在向他宣布这和平之前，和平之子已经在那里，如同他被认识并被预知——被——不是传福音者，而是——天主。因为我们不必害怕会失去它，如果我们无知地向不是和平之子的人传讲，它必归回我们——就是说，那传讲将有益于我们，而非他；但如果所传的和平临到他，就必有益于我们和他两方面。\par

% structure: p0099 source=h2 target=section reason=relative-heading-rank; base=h2

\section{第四十七章——宗徒经文\Quote{祂愿意所有的人都得救。}的另一种解释}

因此，就我们不知道谁会得救而言，天主命令我们愿意所有我们传讲这和平的人得救，并且祂亲自在我们内工作，借着所赐予我们的圣神，将爱倾注在我们心中——也可以这样理解：天主愿意所有的人都得救，因为祂使我们愿意这样；正如\BibleQuote{祂派遣了祂圣子的神，在你们心中呼喊说：阿爸，父啊；}\BibleRef{迦 4:6} 就是说，使我们呼喊：\OriginalTerm{阿爸，父啊}{Abba, Father}。因为论及同一位圣神，祂在另一处说：\BibleQuote{我们领受了义子的神，在祂内我们呼喊：阿爸，父啊！}\BibleRef{罗 8:15} 因此我们呼喊，但那使我们呼喊的被称为呼喊。如果圣经正确地称圣神为呼喊，因为祂使我们呼喊，那么圣经也正确地称天主愿意，因为祂使我们愿意。这样，因为借着责罚，我们只应做两件事：避免离开那与天主相和的平安，或促使那已离开的人重返平安，让我们怀着希望做我们所做的事。如果我们所责罚的人是一个和平之子，我们的平安必临到他；否则，必归回我们。\par

% structure: p0101 source=h2 target=section reason=relative-heading-rank; base=h2

\section{第四十八章——责罚的目的}

因此，尽管有些人的信仰被颠覆，天主的根基却屹立不移，因为主认识属于祂的人，但我们仍不应因此懈怠，忽略责罚那些应受责罚的人。因为并非无故说：\BibleQuote{坏交往败坏好品行；}\BibleRef{格前 15:33} 以及\BibleQuote{那软弱的弟兄，因你的知识而丧亡，基督原为他死了。}\BibleRef{格前 8:11} 我们不要违背这些诫命和有益健康的恐惧，妄自争辩说：\Quote{好吧，让坏交往败坏好品行，让软弱的弟兄丧亡。这与我们何干？天主的根基屹立不移，除了丧亡之子，无人丧亡。} [十六] 愿我们远离这种胡言乱语，以为我们可以在这种疏忽中安然无事。因为确实，除了\OriginalTerm{丧亡之子}{son of perdition}，无人丧亡，但天主借厄则克耳先知的口说：\BibleRef{则 3:18} \BibleQuote{他必要因自己的罪恶而死，但我要向守望者追讨他的血债。}\par

% structure: p0103 source=h2 target=section reason=relative-heading-rank; base=h2

\section{第四十九章——结论}

因此，就我们这些无法分辨哪些是预定得救之人、哪些不是的人来说，我们正为此缘故应当愿意所有人都得救。我们应当以药物般的方式对所有的人施以严厉的责斥，免得他们自己丧亡，或成为毁灭他人的工具。然而，使那责斥对他们有益，是属于天主的权能，祂已预知并预定他们与祂圣子的肖像相符。因为，若我们有时因怕责斥会使人丧亡而避免责斥，那么为什么我们不也责斥，以免人因我们不予责斥而丧亡呢？因为我们所拥有的爱德之心，并不比那位蒙福的宗徒更大，他曾说：\BibleQuote{责斥那些不守规矩的人；安慰那怯懦的人；扶持那软弱的人；对众人要忍耐。要小心，不可有人以恶报恶。}\BibleRef{得前 5:14} 在此应当理解，当一个人应当被责斥却没有被责斥，反而因邪恶的掩饰而被忽略时，便是以恶报恶了。此外，他还说：\BibleQuote{犯罪的人，要在众人面前责斥，使其余的人也有所畏惧。}\BibleRef{弟前 5:20} 这必须理解为是针对那些不隐晦的罪，免得被认为与主的话相违背。因为主说：\BibleQuote{如果你的弟兄得罪了你，你要单独地规劝他，在他和你之间。}\BibleRef{玛 18:15} 然而，主自己将责斥的严厉实行到如此程度，说：\BibleQuote{如果他连教会也不听从，你就把他看作外教人或税吏。}\BibleRef{玛 18:17} 又有谁比祂更爱软弱的人呢？祂为我们众人成了软弱的，并且因那软弱本身为我们众人被钉在十字架上。既然事情如此，恩宠并不阻碍责斥，责斥也不阻碍恩宠；因此，义德当这样被教导，使我们能在忠信祈祷中祈求，借天主的恩宠，所教导的一切得以实行；这两者当以这样的方式并行，使正义的责斥不被忽略。然而，这一切都当以爱德去做，因为爱德既不犯罪，又遮盖许多罪过。\par

\SourceRecord{Translated by Peter Holmes and Robert Ernest Wallis, and revised by Benjamin B. Warfield.；Nicene and Post-Nicene Fathers, First Series；Vol. 5.；Edited by Philip Schaff.；Buffalo, NY: Christian Literature Publishing Co.,；1887.；<http://www.newadvent.org/fathers/1513.htm>.}
